\documentclass[10pt,aspectratio=169]{beamer}
% Preámbulo modular para presentaciones Beamer (Modelación Estadística)
% Diseño y paleta institucional del Tecnológico de Monterrey

\documentclass[aspectratio=169,xcolor={dvipsnames,table}]{beamer}

\usepackage[utf8]{inputenc}
\usepackage[spanish,mexico]{babel}
\usepackage{amsmath,amssymb,amsfonts,mathtools}
\usepackage{graphicx}
\usepackage{listings}
\usepackage{booktabs}
\usepackage{tikz}
\usetikzlibrary{matrix,arrows,decorations.pathmorphing,shapes.geometric,positioning}

% Paleta Institucional Tecnológico de Monterrey
\definecolor{TecAzul}{HTML}{0039A6}       % Azul TEC: color primario institucional
\definecolor{TecAzulOscuro}{HTML}{1A2E51} % Azul marino oscuro
\definecolor{TecRojo}{HTML}{EC2661}       % Rojo de acento (paleta miTec)
\definecolor{TecGris}{HTML}{646464}       % Gris neutro para comentarios y subtítulos
\definecolor{TecFondoBloque}{HTML}{F4F6F9}% Fondo suave para bloques

% Configuración de Tema Beamer
\usetheme{Boadilla}
\usecolortheme[named=TecAzulOscuro]{structure}
\setbeamercolor{alerted text}{fg=TecRojo}
\setbeamercolor{palette primary}{bg=TecAzulOscuro,fg=white}
\setbeamercolor{palette secondary}{bg=TecAzul,fg=white}
\setbeamercolor{palette tertiary}{bg=TecAzulOscuro!80!black,fg=white}
\setbeamercolor{title}{fg=TecAzulOscuro,bg=TecFondoBloque}
\setbeamercolor{frametitle}{fg=TecAzulOscuro,bg=TecFondoBloque}

% Bloques personalizados elegantes
\setbeamercolor{block title}{bg=TecAzulOscuro,fg=white}
\setbeamercolor{block body}{bg=TecFondoBloque,fg=black}
\setbeamercolor{block title alerted}{bg=TecRojo,fg=white}
\setbeamercolor{block body alerted}{bg=TecRojo!8,fg=black}
\setbeamercolor{block title example}{bg=TecAzul,fg=white}
\setbeamercolor{block body example}{bg=TecAzul!8,fg=black}

% Redondear bordes y añadir sombra sutil
\setbeamertemplate{blocks}[rounded][shadow=true]
\setbeamertemplate{navigation symbols}{}

% Pie de página limpio con número de diapositiva
\setbeamertemplate{footline}{
  \leavevmode%
  \hbox{%
  \begin{beamercolorbox}[wd=.7\paperwidth,ht=2.25ex,dp=1ex,left]{author in head/foot}%
    \hspace*{2ex}\insertshortauthor~~(\insertshortinstitute) \hfill \insertshorttitle
  \end{beamercolorbox}%
  \begin{beamercolorbox}[wd=.3\paperwidth,ht=2.25ex,dp=1ex,right]{date in head/foot}%
    \insertshortdate{}\hspace*{2em}
    \insertframenumber{} / \inserttotalframenumber\hspace*{2ex}
  \end{beamercolorbox}}%
  \vskip0pt%
}

% Configuración de Listings de Python para visualización pedagógica de código
\lstset{
    language=Python,
    basicstyle=\ttfamily\footnotesize,
    keywordstyle=\color{TecAzulOscuro}\bfseries,
    stringstyle=\color{TecRojo},
    commentstyle=\color{TecGris}\itshape,
    showstringspaces=false,
    breaklines=true,
    frame=lines,
    rulecolor=\color{TecAzulOscuro!30},
    backgroundcolor=\color{TecFondoBloque},
    aboveskip=0.5em,
    belowskip=0.5em,
    literate={á}{{\'a}}1 {é}{{\'e}}1 {í}{{\'i}}1 {ó}{{\'o}}1 {ú}{{\'u}}1
             {Á}{{\'A}}1 {É}{{\'E}}1 {Í}{{\'I}}1 {Ó}{{\'O}}1 {Ú}{{\'U}}1
             {ñ}{{\~n}}1 {Ñ}{{\~N}}1 {¿}{{?`}}1 {¡}{{!`}}1
}

% Metadatos institucionales por defecto
\title[Modelación Estadística]{Modelación Estadística para Ciencia de Datos e Ingeniería}
\author[J. Castillo Colmenares]{Juliho David Castillo Colmenares}
\institute[Tec de Monterrey]{Tecnológico de Monterrey \\ Escuela de Ingeniería y Ciencias}
\date{\today}

% Comandos y macros estadísticas compartidas para las presentaciones Beamer (Metropolis)

% Conjuntos numéricos básicos
\newcommand{\R}{\mathbb{R}}
\newcommand{\N}{\mathbb{N}}
\newcommand{\Q}{\mathbb{Q}}
\newcommand{\Z}{\mathbb{Z}}
\newcommand{\set}[1]{\left\{ #1 \right\}}
\newcommand{\abs}[1]{\left|#1\right|}
\newcommand{\norm}[1]{\left\| #1 \right\|}

% Comandos estadísticos y probabilísticos del curso
\newcommand{\Var}{\operatorname{Var}}
\newcommand{\cov}{\operatorname{Cov}}
\newcommand{\corr}{\rho}
\newcommand{\comb}[2]{\begin{pmatrix} #1 \\ #2 \end{pmatrix}}
\newcommand{\s}{\sigma}
\newcommand{\particion}{\mathfrak{P}}
\newcommand{\card}[1]{\#\left( #1 \right)}
\newcommand{\seq}[3]{\set{#1}_{#2}^{#3}}
\newcommand{\E}{\mathbb{E}}
\newcommand{\Prob}{\mathbb{P}}

% Entornos pedagógicos y cajas en español académico natural
\newenvironment{discusion}{%
  \begin{alertblock}{Pregunta de Discusión}%
}{%
  \end{alertblock}%
}

\newenvironment{reflexion}{%
  \begin{alertblock}{Reflexión para el Aula}%
}{%
  \end{alertblock}%
}

\newenvironment{paradoja}{%
  \begin{alertblock}{Paradoja de Exploración}%
}{%
  \end{alertblock}%
}

\newenvironment{motivacion}{%
  \begin{exampleblock}{Motivación: Ciencia de Datos en la Práctica}%
}{%
  \end{exampleblock}%
}

\newenvironment{retoclase}{%
  \begin{block}{Reto Rápido en Equipo}%
}{%
  \end{block}%
}


\title[Distribuciones de funciones]{Sección 5.2: Distribuciones de probabilidad de funciones de variable aleatoria}
\subtitle{Preimágenes, superficies de nivel y sumas}
\author[J. Castillo Colmenares]{Juliho Castillo Colmenares}
\institute{Tecnológico de Monterrey}
\date{}

\begin{document}
\begin{frame}[plain]\vspace{-0.8cm}\titlepage\end{frame}

\begin{frame}{Hoja de ruta --- capítulo 5}
  \footnotesize
  \begin{itemize}\itemsep=0.08cm
    \item \textbf{05.00} Introducción inferencial
    \item \textbf{05.01} Transformación de variables
    \item \textbf{05.02} Funciones de variables aleatorias \emph{(hoy)}
    \item \textbf{05.03--05.05} Distribuciones muestrales clásicas
  \end{itemize}
  \begin{block}{Objetivo didáctico}
    Generalizar el cambio de variable a funciones no monótonas y a varias
    variables, y reconocer la distribución Erlang de una suma de exponenciales.
  \end{block}
\end{frame}

\begin{frame}{Motivación: una función cambia la variable}
  \footnotesize
  Cuando $Y=g(X)$, la densidad de $Y$ debe reunir toda la probabilidad de los
  valores de $X$ que producen el mismo resultado.
  \begin{alertblock}{Pregunta guía}
    ¿Cuántas preimágenes tiene cada $y$ y cómo cambia la escala local?
  \end{alertblock}
\end{frame}

\begin{frame}{Una variable: fórmula general}
  \footnotesize
  Para una función diferenciable con preimágenes regulares:
  \[
    f_Y(y)=\sum_{x:g(x)=y}f_X(x)\left|\frac{dx}{dy}\right|
    =\sum_{x:g(x)=y}\frac{f_X(x)}{|g'(x)|}.
  \]
  La suma recorre todas las soluciones de $g(x)=y$.
\end{frame}

\begin{frame}{Condición de regularidad}
  \footnotesize
  \begin{block}{Hipótesis local}
    La condición $|g'(x)|\neq0$ en cada preimagen garantiza que la función es
    localmente invertible en ese punto.
  \end{block}
  \begin{exampleblock}{Consecuencia}
    Si la derivada se anula, la fórmula requiere tratar con cuidado el punto
    crítico y el soporte resultante.
  \end{exampleblock}
\end{frame}

\begin{frame}{Varias variables: superficie de nivel}
  \footnotesize
  Si $Y=g(X_1,\ldots,X_n)$ tiene densidad conjunta,
  \[
    f_Y(y)=\int_{g(x_1,\ldots,x_n)=y}
    f_{X_1,\ldots,X_n}(x_1,\ldots,x_n)\frac{dS}{|\nabla g|}.
  \]
  \begin{block}{Geometría}
    La integral recorre la superficie de nivel que reúne todas las
    configuraciones con el mismo valor de $Y$.
  \end{block}
\end{frame}

\begin{frame}{Independencia y sumas}
  \footnotesize
  Para $Y=X_1+X_2$, la superficie de nivel es la recta
  $x_1+x_2=y$. Si $X_1,X_2$ son independientes, la densidad conjunta es el
  producto de sus densidades.
  \begin{alertblock}{Puente}
    La integración sobre esa recta produce la densidad de la suma.
  \end{alertblock}
\end{frame}

\begin{frame}{Lectura de la sección}
  \footnotesize
  \begin{itemize}\itemsep=0.12cm
    \item Una variable: suma sobre preimágenes.
    \item Varias variables: integral sobre una superficie de nivel.
    \item El gradiente $\nabla g$ mide el cambio local de escala.
    \item La independencia simplifica la densidad conjunta.
  \end{itemize}
\end{frame}

\begin{frame}{Aplicación derivada 1/4: localizar preimágenes}
  \footnotesize
  Para cada $y$, se resuelve $g(x)=y$ antes de evaluar la fórmula.
  \begin{block}{Procedimiento}
    \begin{enumerate}\itemsep=0.1cm
      \item hallar todas las soluciones $x$;
      \item verificar $g'(x)\neq0$;
      \item sumar las contribuciones.
    \end{enumerate}
  \end{block}
\end{frame}

\begin{frame}{Aplicación derivada 1/4: interpretación}
  \footnotesize
  \[
    \frac{f_X(x)}{|g'(x)|}
  \]
  es la contribución de una sola rama. La densidad de $Y$ es la suma de todas
  las ramas que llegan al mismo valor.
\end{frame}

\begin{frame}{Aplicación derivada 2/4: puntos críticos}
  \footnotesize
  Si $g'(x)=0$, el factor $1/|g'(x)|$ no puede evaluarse directamente.
  \begin{alertblock}{Diagnóstico}
    Los puntos críticos suelen marcar extremos del soporte o singularidades de
    la densidad; deben analizarse con el límite y la integral total.
  \end{alertblock}
\end{frame}

\begin{frame}{Aplicación derivada 2/4: interpretación}
  \footnotesize
  La condición de regularidad no es un detalle formal: evita contar una
  transformación como si fuera invertible donde no lo es.
  \begin{exampleblock}{Regla de lectura}
    Primero se identifica la geometría de $g$; después se aplica la fórmula.
  \end{exampleblock}
\end{frame}

\begin{frame}{Aplicación resuelta 3/4: suma exponencial}
  \footnotesize
  Sean $X_1,X_2\sim\operatorname{Exp}(1)$ independientes y $Y=X_1+X_2$.
  Para $y>0$, la recta de nivel tiene $0<x_1<y$ y $x_2=y-x_1$.
  \[
    f_{X_1,X_2}(x_1,x_2)=e^{-(x_1+x_2)}.
  \]
\end{frame}

\begin{frame}{Aplicación resuelta 3/4: integración}
  \footnotesize
  \[
    f_Y(y)=\int_0^y e^{-(x_1+y-x_1)}\,dx_1
    =\int_0^y e^{-y}\,dx_1=ye^{-y},\quad y>0.
  \]
  \begin{alertblock}{Resultado}
    La suma tiene densidad Erlang de forma $2$, equivalente a
    $\operatorname{Gamma}(2,1)$.
  \end{alertblock}
\end{frame}

\begin{frame}{Aplicación resuelta 4/4: generalización}
  \footnotesize
  La sección concluye que la suma de $n$ exponenciales iid tiene distribución
  \[
    \operatorname{Gamma}(n,1).
  \]
  \begin{block}{Lectura}
    El caso de dos variables muestra cómo una integral sobre una superficie de
    nivel produce una familia conocida.
  \end{block}
\end{frame}

\begin{frame}{Aplicación resuelta 4/4: soporte}
  \footnotesize
  Para la suma exponencial, $X_1>0$ y $X_2>0$ implican $Y>0$.
  \begin{exampleblock}{Comprobación}
    La densidad $ye^{-y}$ se interpreta únicamente para $y>0$ y su integral en
    ese soporte es uno.
  \end{exampleblock}
\end{frame}

\begin{frame}{Puente computacional 1/4: mapa geométrico}
  \footnotesize
  \begin{tabular}{ll}
    \hline
    Transformación & Conjunto que se integra \\
    \hline
    $Y=g(X)$ & preimágenes $g^{-1}(y)$ \\
    $Y=g(X_1,X_2)$ & curva $g(x_1,x_2)=y$ \\
    $Y=g(X_1,\ldots,X_n)$ & superficie de nivel \\
    \hline
  \end{tabular}
\end{frame}

\begin{frame}{Puente computacional 2/4: densidad conjunta}
  \footnotesize
  Para variables independientes,
  \[
    f_{X_1,X_2}(x_1,x_2)=f_{X_1}(x_1)f_{X_2}(x_2).
  \]
  Esta factorización es la entrada de la integral sobre la superficie de nivel.
\end{frame}

\begin{frame}{Puente computacional 3/4: gradiente}
  \footnotesize
  El término $|\nabla g|$ ajusta el elemento de superficie:
  \[
    dS/|\nabla g|.
  \]
  \begin{block}{Interpretación}
    El jacobiano generalizado convierte una medida geométrica en masa de
    probabilidad.
  \end{block}
\end{frame}

\begin{frame}{Puente computacional 4/4: lista de verificación}
  \footnotesize
  \begin{enumerate}\itemsep=0.12cm
    \item fijar el soporte de $Y$;
    \item describir todas las preimágenes o niveles;
    \item introducir la densidad conjunta;
    \item integrar y comprobar la normalización.
  \end{enumerate}
\end{frame}

\begin{frame}{Síntesis de la sección}
  \footnotesize
  \begin{itemize}\itemsep=0.12cm
    \item Las preimágenes se suman en una dimensión.
    \item Las superficies de nivel se integran en varias dimensiones.
    \item El gradiente controla el cambio de escala.
    \item Sumas de exponenciales producen una familia Gamma/Erlang.
  \end{itemize}
\end{frame}

\begin{frame}{Transición curricular}
  \footnotesize
  \begin{block}{Lo que queda establecido}
    La distribución de una función se obtiene reuniendo todos los puntos de la
    muestra que generan el mismo valor transformado.
  \end{block}
  \begin{alertblock}{Siguiente sección}
    \textbf{05.03 Distribución chi-cuadrada}: una distribución muestral clásica
    construida a partir de cuadrados de variables normales.
  \end{alertblock}
\end{frame}
\end{document}
